\chapter*{Заключение}
\addcontentsline{toc}{chapter}{Заключение}
\label{ch:conclusion}

В ходе выполнения лабораторной работы~\cfgLabNumber{} было развёрнуто и
настроено облачное файловое хранилище \textbf{Seafile} на виртуальной машине
под управлением Ubuntu~22.04~LTS.

Выполненные задачи:
\begin{itemize}
  \item создана ВМ \texttt{seafile} с Ubuntu~22.04, назначен статический
        IP \texttt{192.168.\cfgLabStudentN.4/24};
  \item хост зарегистрирован в DNS-зоне через запись~A на сервере \texttt{gateway};
  \item установлены Python~3, системные зависимости и pip-пакеты
        (Django, Pillow, mysqlclient, pycryptodome~и~др.);
  \item развёрнута и настроена \textbf{MariaDB}: установлен пароль root,
        автозапуск активирован;
  \item загружен дистрибутив \texttt{seafile-server\_9.0.9},
        выполнен скрипт \texttt{setup-seafile-mysql.sh},
        созданы три базы данных (ccnet-db, seafile-db, seahub-db);
  \item настроен обратный прокси \textbf{Nginx}: запросы на
        \texttt{192.168.\cfgLabStudentN.4:80} проксируются в Seahub
        (\texttt{127.0.0.1:8000});
  \item создан \texttt{seafile.service} для автозапуска через \texttt{systemd};
  \item через веб-консоль создан аккаунт администратора и пользователь;
  \item на ВМ Desktop установлен клиент \texttt{seafile-gui},
        произведена синхронизация библиотеки с сервером.
\end{itemize}

Полученные навыки — развёртывание корпоративного облачного хранилища,
работа с MariaDB, настройка Nginx как обратного прокси и управление
systemd-сервисами — применяются при построении корпоративной IT-инфраструктуры.

\endinput